% Part: incompleteness
% Chapter: theories-computability
% Section: inseparability

\documentclass[../../../include/open-logic-section]{subfiles}

\begin{document}

\olfileid[hi]{inc}{tcp}{ins}

\olsection{$\Th{Q}$ में सिद्ध और खंडनीय वाक्य संगणनीयतः अवियोज्य हैं}

$\Th{\bar Q}$ को उन वाक्यों का समुच्चय मानें जिनके {\em निषेध}
$\Th{Q}$ में सिद्ध होते हैं, अर्थात्
$\Th{\bar Q} = \Setabs{!A}{\Th{Q} \Proves \lnot !A}$। याद करें कि
असंयुक्त समुच्चय $A$ और $B$ संगणनीयतः अवियोज्य कहलाते हैं यदि ऐसा कोई
संगणनीय समुच्चय $C$ न हो कि $A \subseteq C$ और
$B \subseteq \Complement{C}$।

\begin{lem}
$\Th{Q}$ और $\Th{\bar Q}$ संगणनीयतः अवियोज्य हैं।
\end{lem}

\begin{proof}
मान लें $C$ ऐसा संगणनीय समुच्चय है कि $\Th{Q} \subseteq C$ और
$\Th{\bar Q} \subseteq \Complement{C}$। $R(x,y)$ को संबंध
\begin{quote}
$x$ किसी सूत्र $!D(u)$ का कूट है और $!D(\num y)$,~$C$ में है
\end{quote}
मानें। हम दिखाएँगे कि $R(x,y)$ कोई सार्वत्रिक संगणनीय संबंध है, जिससे
विरोधाभास मिलेगा।

मान लें $S(y)$ संगणनीय है और~$\Th{Q}$ में $!D_S(u)$ द्वारा निरूपित है। तब
\begin{eqnarray*}
S(\num n) & \lif & \Th{Q} \Proves !D_S(\num n) \\
& \lif & {}!D_S(\num n) \in C
\end{eqnarray*}
और
\begin{eqnarray*}
\lnot S(\num n) & \lif & \Th{Q} \Proves \lnot !D_S(\num n) \\
& \lif & {}!D_S(\num n) \in \Th{\bar Q} \\
& \lif & {}!D_S(\num n) \not\in C।
\end{eqnarray*}
अतः $S(y)$,~$R(\#(!D_S(\num u)),y)$ के समतुल्य है।
\end{proof}

\end{document}
